\section{Open Questions}

The following five questions extend the paper's central mechanistic
question --- \emph{how does dose deposition drive alveolar-scale
fibrosis?} --- into directions the paper does not itself address.
They are chosen to be genuinely open (no straightforward answer in
the current literature), concrete (each has a runnable next step),
and free-endpoint (no paywalled data, no wet lab).

\begin{enumerate}
  \item \textbf{FLASH regime extension.} Does the ABM-MC coupled model,
    extended to FLASH-radiotherapy dose rates ($\sim 40$~Gy/s), predict
    reduced fibrosis at matched total dose vs.\ conventional dose-rate,
    and does that match preclinical FLASH lung-sparing data?
    \emph{Next steps:} extend the LQ kernel with a dose-rate term
    (temporal integration coupled to a transient O$_2$ field);
    recompile and run 1-fx conventional vs.\ FLASH at 10, 20, 30~Gy,
    10 reps; compare late $\Delta$ECM(D) and FSU survival vs.\ Vozenin
    et al.\ and Favaudon et al.\ mouse-lung FLASH datasets.

  \item \textbf{Omics-informed cell-state parameterisation.} Can the
    ABM's cell-state definitions be re-parameterised from single-cell
    RNA-seq atlases of irradiated (or bleomycin-injured) lung, and do
    the omics-informed transition rates change the predicted
    $\Delta$ECM(D) shape?
    \emph{Next steps:} download Strunz et al.\ GSE141259; fit pseudotime
    transition rates for AEC2$\rightarrow$damaged$\rightarrow$
    senescent$\rightarrow$apoptotic via CellRank or Monocle3; replace
    ABM constants with pseudotime-derived rates; re-run dose sweep;
    compare ED$_{50}$ shift.

  \item \textbf{Clinical cohort mapping.} How well does the model's
    predicted $\Delta$ECM(D) map to specific clinical cohorts (stage III
    NSCLC post-SBRT, thoracic lymphoma post-mantle), and does the
    fibrosis-index predict grade-2+ pneumonitis in the QUANTEC
    datasets?
    \emph{Next steps:} extract MLD and V20 from QUANTEC (Marks et al.\
    2010 Table 2); run ABM at each cohort's MLD as equivalent uniform
    dose; compute RSI at 6 and 12 months; ROC against pneumonitis
    grade-2+ rates; compare AUC to Lyman-Kutcher-Burman NTCP.

  \item \textbf{AEC2 cell-cycle sensitivity.} How sensitive is the
    predicted fibrosis-index to AEC2 cell-cycle length assumptions,
    and could this parameter explain between-patient variance in RILF
    risk?
    \emph{Next steps:} parameterise the repopulation timescale
    $\tau_{\mathrm{prol}} \in \{2,4,7,14\}$ days; run 1-fx sweep at each
    $\tau$, 10 reps; compute CV of $\Delta$ECM across $\tau$;
    compare to clinical between-patient RILF variance from open cohort
    data.

  \item \textbf{ML-driven per-patient personalisation from CT.} Can
    machine-learning parameter identification from serial post-RT CT
    (via radiomics-derived ECM-density surrogates) recover ABM
    parameters, enabling per-patient personalisation?
    \emph{Next steps:} build a neural surrogate of the ABM mapping
    $\{\alpha,\beta,\mathrm{phag\_fraction},k_{\mathrm{mf\_grow}}\}$ to
    $\Delta$ECM(t); train on 10k Latin-hypercube ABM runs; download
    TCIA NSCLC-Radiomics collection; extract time-series HU-density
    change in treated volume; invert the surrogate for patient-specific
    parameters; test whether parameters cluster by pneumonitis grade.
\end{enumerate}
